\documentclass[12pt,a4paper,UTF8]{ctexart}

% 必要宏包
\usepackage{geometry}
\usepackage{amsmath}
\usepackage{amssymb}
\usepackage{graphicx}
\usepackage{booktabs}
\usepackage{hyperref}
\usepackage{float}
\usepackage{enumitem}
\usepackage{multirow}

% 页面布局
\geometry{left=3cm,right=2.5cm,top=2.5cm,bottom=2.5cm}

% 超链接样式
\hypersetup{
    colorlinks=true,
    linkcolor=blue,
    citecolor=blue,
    urlcolor=blue
}

% 论文信息
\title{\textbf{基于耦合多智能体强化学习的\\低空航空交通协同管控模型}}
\author{作者：彭义森 \quad 指导老师：韩云祥 \\ 四川大学}
\date{2026年7月}

\begin{document}

\maketitle

\begin{abstract}
随着先进空中交通（Advanced Air Mobility, AAM）的快速发展，电动垂直起降飞行器（eVTOL）在城市低空环境中的高密度运行，对传统空中交通管理系统提出了严峻挑战。现有多智能体强化学习（MARL）方法普遍存在三大核心缺陷：一是采用静态图结构或全连接网络编码智能体关联，无法捕捉飞行器间动态时空耦合关系；二是训练过程中冲突样本高度稀疏导致学习效率低下；三是对低空环境的物理约束与不确定性建模不足。针对城市低空十字交叉路口场景下的多飞行器冲突避免与流量协同优化问题，本文提出了一种基于耦合多智能体强化学习的低空交通管控模型。该模型通过动态图注意力网络（DyGAT）同时融合空间注意力、时序注意力与动态边权重计算，刻画飞行器间的时空耦合关联；设计了分层奖励函数兼顾飞行安全与运行效率；同时引入课程学习与基于冲突事件的轨迹加权机制解决稀疏奖励与收敛稳定性问题。在融合风扰与传感器噪声的仿真环境中，所提模型在21架eVTOL的交叉路口场景中取得了最优性能：相较于传统固定时序调度（RR）方法，冲突次数降低57.2\%，平均运行延迟缩短63.9\%；相较于当前最优的MARL方法MAPPO，冲突次数降低10.3\%，平均运行延迟缩短12.3\%，累计奖励提升44.2\%。本文工作为城市低空交通的分布式自主管控提供了一种有潜力的技术路线。

\noindent \textbf{关键词：} 低空航空交通管理；多智能体强化学习；动态图注意力网络；冲突避免；集中训练分布式执行；时空耦合建模
\end{abstract}

%==========================================================================
\section{引言}
%==========================================================================

城市地面交通拥堵问题的日益加剧，推动了以eVTOL为核心的城市空中交通（Urban Air Mobility, UAM）产业的快速发展。美国联邦航空管理局（FAA）预测，到2040年将有数百万架无人机与eVTOL飞行器在美国低空空域运行。与传统民航航路的结构化交通不同，城市低空交通具有高密度、强动态、高异构性的特征，飞行器在交叉路口与汇合点的协同管控与冲突避免，已成为低空无人机交通管理系统（Unmanned Aircraft System Traffic Management, UTM）面临的核心挑战。

传统空中交通流量管理（Air Traffic Flow Management, ATFM）方法多采用集中式优化架构（如整数规划、蒙特卡洛仿真等），依赖预先规划的飞行计划与静态航路设计，难以适配低空环境中天气扰动、传感器噪声与突发流量变化等动态不确定性因素。近年来，多智能体强化学习（Multi-Agent Reinforcement Learning, MARL）因其分布式决策与自适应学习的特性，被成功引入交通管控领域。Agogino与Tumer系列工作率先为空中交通流量管理建立了多智能体学习范式：从分布式航路点智能体调控框架\cite{ref1}，到基于差异奖励的大规模空中交通流量管理方法\cite{ref2}，系统性地验证了MARL在ATM领域的有效性。此后，Brittain与Wei针对AAM场景中的高密度航路段汇合与交叉路口问题，设计了基于A2C结合PPO损失函数的深度MARL框架，在BlueSky仿真器中实现了99.97\%的冲突规避率\cite{ref3}。Tumer与Agogino进一步为空中交通流量管理提出了基于分布式智能体的反馈控制机制\cite{ref4}，并在IEEE Intelligent Systems上综述了学习型多智能体系统提升ATM性能的关键技术路径\cite{ref5}。Xie等将强化学习与遗传算法相结合，构建了面向城市空中交通需求-容量平衡的混合学习框架\cite{ref6}。

尽管上述工作取得了显著进展，现有研究仍存在三个关键且未被充分解决的研究空白：

\subsection{研究空白一：时空耦合关系刻画不足}

现有MARL方法大多通过全连接网络或简单静态图结构编码智能体间的关联，无法捕捉飞行器间动态变化的邻域交互关系。具体表现为：(1)~Brittain与Wei虽然采用了分布式架构，但仅考虑最近邻飞行器的静态关联，当飞行器快速移动时邻域拓扑结构发生剧烈变化，静态图无法及时更新关联权重\cite{ref3}；(2)~Spatharis等提出的层级式MARL方案虽将状态-动作空间进行抽象以提升搜索效率，但智能体间的交互仍通过预定义关系矩阵编码，缺乏对空间拓扑结构的显式建模\cite{ref7}；(3)~如张春阳与席雅琪在综述中所指出\cite{ref8}，Wu等提出的MARDDPG结合LSTM增强了模型适应性，但仅通过全连接网络融合邻域信息，无法按距离、速度与航向角区分不同邻域飞行器的冲突风险差异。Ghosh等提出的深度集成MARL方法通过元策略仲裁不同专家策略，在高密度场景中表现出色，但其状态编码仍依赖网格化的密度图（Geo-hash），缺少对飞行器间细粒度时空耦合关系的显式建模\cite{ref9}。这些缺陷导致现有模型在高密度场景下冲突规避能力急剧下降——当飞行器数量超过15架时，冲突率呈指数级上升。

\subsection{研究空白二：稀疏奖励问题未得到有效解决}

低空交通中冲突事件属于低概率高风险事件，在训练过程中冲突样本占比通常低于0.1\%，导致有效训练信号极度稀疏。Aslani等在城市地面交通信号控制的研究中指出，标准的$\epsilon$-greedy探索策略在高维连续动作空间中难以高效覆盖稀疏的关键状态区域\cite{ref10}；翟子洋等的综述系统梳理了大规模交通信号控制中RL与DRL方法的发展脉络，其中明确指出稀疏奖励是MARL在多智能体交通场景中面临的首要挑战\cite{ref11}。具体而言：(1)~标准经验回放机制中冲突样本被采样的概率极低，模型需要数百万次交互才能学习到基本的冲突规避策略；(2)~冲突规避的奖励信号通常只有在冲突发生或成功通过交叉路口时才会出现，导致信用分配困难；(3)~随机探索策略难以覆盖罕见但致命的极端冲突构型。Ghavamzadeh等提出的层级式MARL框架虽然通过任务分解降低了学习复杂度\cite{ref12}，但层级结构的先验设计依赖领域知识，难以直接泛化到动态变化的飞行器密度场景。

\subsection{研究空白三：物理约束与不确定性建模粗糙}

多数研究采用简化的环境模型，忽略了低空环境的复杂性与不确定性。Spatharis等已指出，传统固定安全距离无法适配不同飞行速度下的制动需求——速度越快所需安全裕度越大，且真实低空环境中的风扰与传感器噪声对飞行安全有重大影响\cite{ref7}。然而，大多数现有研究将飞行器视为质点，忽略了最大加速度、最大转弯角等物理运动约束，导致生成的控制指令无法被真实飞行器执行。

\subsection{本文研究目标与核心贡献}

针对上述三个关键研究空白，本文的核心研究目标是：\textbf{设计一种能够精准刻画飞行器间动态时空耦合关系、高效利用稀疏冲突样本、且充分考虑低空环境物理约束与不确定性的多智能体强化学习管控模型}。

本文的主要贡献包括：

\textbf{贡献一：面向低空交叉路口冲突规避的仿真环境构建。} 构建了融合风扰与传感器噪声的仿真环境\texttt{AirTrafficEnv}，引入与速度相关的动态安全距离模型，严格遵循eVTOL的物理运动约束（最大加速度3.0~m/s$^2$、最大转弯角$\pi/6$等）。

\textbf{贡献二：基于DyGAT的耦合智能体群网络。} 提出基于动态图注意力网络（DyGAT）的耦合智能体群架构，通过距离、速度、航向角三个维度的动态边权重计算实时刻画飞行器间关联强度，同时融合空间注意力与时序注意力。消融实验表明该模块使冲突规避能力提升174\%。

\textbf{贡献三：课程学习与冲突事件轨迹加权训练框架。} 设计渐进式课程学习机制（5架逐步递增至21架），并在PPO框架内引入基于冲突事件的轨迹加权策略以提升稀疏样本利用效率。

\textbf{贡献四：全面系统的实验验证与性能基准。} 与8种基线方法进行对比（涵盖规则方法、单智能体RL、独立多智能体学习与CTDE-MARL四类），通过消融实验、鲁棒性分析与敏感性分析建立可复现的性能基准。

本文其余部分结构如下：第2节系统回顾相关工作并定位本文方法的差异性；第3节进行问题建模并详细阐述管控模型的设计细节；第4节展示实验结果与深入分析；第5节进行讨论并总结全文。

%==========================================================================
\section{相关工作}
%==========================================================================

本节从传统空中交通流量管理方法、强化学习在交通管控中的应用两个维度系统梳理相关研究，并在表\ref{tab:related_work}中给出本文与代表性工作的系统性对比。

\subsection{传统空中交通流量管理方法}

传统ATFM方法可分为集中式优化与分布式启发式两类。集中式优化方面，Bertsimas等采用整数规划求解空域流量分配问题，Menon等通过交通流欧拉模型的线性规划实现动态调控\cite{ref6}。这类方法在民航干线等可预测场景中效果良好，但在低空高密度、强动态场景中存在计算复杂度高与响应滞后的问题。启发式方法方面，Xie等通过遗传算法优化无人机4D航迹以求解需求-容量平衡问题\cite{ref6}；Agogino与Tumer设计了基于规则的反馈控制机制，通过交通流量反馈动态调整航路点放行间隔，相比固定轮询调度显著降低了队列长度\cite{ref4}。此类方法计算效率高，但依赖人工设计的规则与启发式函数，泛化能力受限。

\subsection{强化学习在交通管控中的应用}

在单智能体强化学习方面，深度Q网络（DQN）已被应用于单飞行器冲突规避与单交叉口信号控制等场景。王迎在其硕士论文中构建了"预测-决策"协同的深度强化学习城市交通信号控制框架，验证了LSTM与DDQN在SUMO仿真平台上的有效性\cite{ref13}。赵晓华等将Q-learning与BP神经网络结合，在PARAMICS仿真环境中实现了交叉口信号灯的自适应控制\cite{ref14}。承向军等提出了基于Q-学习的单路口交通信号控制方法，通过模糊控制规则集与Q-学习优化相结合，使平均停车延迟相比定时控制降低了29.25\%\cite{ref15}。赖建辉设计了基于D3QN的深度强化学习交通控制优化方法，使用离散交通状态编码与卷积拼接的路网状态表示，在杭州真实路网数据上验证了算法的工程可行性\cite{ref16}。Aslani等在城市地面交通信号控制中系统对比了DQN与A2C在不同交通扰动事件下的鲁棒性，验证了Actor-Critic架构在连续动作空间中的优势\cite{ref10}。张新武在其硕士论文中提出了基于PPO系列算法的交通信号控制优化方案\cite{ref17}，在单路口与多路口场景中均验证了PPO相比Q-learning与Double DQN的优越性。卢守峰等以平均排队长度差最小为优化目标，构建了单交叉口在线Q学习模型\cite{ref18}。江波等将深度Q学习应用于机场道口的协同智能调速，在双向多航道场景中实现了99\%以上的通行成功率\cite{ref19}。

在多智能体强化学习领域，Ghavamzadeh等提出了层级式MARL框架\cite{ref12}，通过将协作学习分解为任务层级减少了状态-动作空间的探索负担，为后续层级式ATM方法提供了理论基础。Tumer与Agogino的工作\cite{ref2,ref4,ref5}系统性地建立了基于航路点智能体的分布式ATM学习框架，通过差异奖励函数$D_i = G(z) - G(z_{-i})$有效解决了信用分配问题。Ghosh等提出了深度集成MARL方法\cite{ref9}，通过元策略网络在基于核的模型驱动方法与基于PPO的深度学习方法之间进行动态仲裁。翟子洋等的综述\cite{ref11}全面梳理了大规模交通信号控制中从独立Q-learning到QMIX、MAPPO等CTDE方法的演进路线，并指出图神经网络（GAT、GCN）在刻画路网拓扑关系方面的前沿应用。

\subsection{本文方法与现有工作的差异化定位}

表\ref{tab:related_work}将本文方法与上述代表性工作进行系统性对比。从表中可以看出：现有工作在时空耦合建模方面多采用全连接网络或静态图结构，未能同时融合空间注意力、时序注意力与动态边权重；在训练机制方面多依赖标准经验回放或基础PPO框架，未针对冲突样本稀疏问题进行专门设计；在环境建模方面多采用2D简化模型，未充分考虑物理约束与扰动。本文方法在这些维度上均做出了针对性改进。

\begin{table}[H]
    \centering
    \caption{本文方法与代表性相关工作的系统性对比}
    \label{tab:related_work}
    \begin{tabular}{|c|c|c|c|c|c|}
        \hline
        \textbf{方法} & \textbf{时空耦合建模} & \textbf{图结构类型} & \textbf{稀疏奖励处理} & \textbf{物理约束} & \textbf{扰动建模} \\
        \hline
        Brittain \& Wei \cite{ref3} & LSTM+FC & 静态全连接 & $\times$ & 部分 & $\times$ \\
        \hline
        Ghosh et al. \cite{ref9} & Geo-hash密度图 & 网格化 & $\times$ & 部分 & $\times$ \\
        \hline
        Spatharis et al. \cite{ref7} & 预定义关系矩阵 & 层级分组 & $\times$ & $\times$ & $\times$ \\
        \hline
        QMIX/MAPPO基线 & FC/RNN & 全连接 & $\times$ & 部分 & $\times$ \\
        \hline
        \textbf{本文方法} & \textbf{DyGAT(空间+时序+动态边)} & \textbf{动态图} & \textbf{课程学习+轨迹加权} & \textbf{完整} & \textbf{风扰+噪声} \\
        \hline
    \end{tabular}
\end{table}

%==========================================================================
\section{问题建模与管控模型设计}
%==========================================================================

本文聚焦于城市低空的十字交叉路口场景：$N$架eVTOL从不同方向飞向交叉中心，需在保持安全间隔的前提下，以最小的延迟通过交叉区域。本节首先将该问题建模为分布式马尔可夫决策过程（Dec-MDP），然后依次阐述环境模型、奖励函数、DyGAT网络、单智能体网络、GraphMixer混合网络与训练优化机制。

\subsection{分布式马尔可夫决策过程建模}

将$N$架eVTOL的协同管控问题形式化定义为Dec-MDP，由六元组$\langle \mathcal{N}, \mathcal{S}, \mathcal{A}, \mathcal{P}, \mathcal{R}, \gamma \rangle$描述：

\begin{itemize}
    \item \textbf{智能体集合} $\mathcal{N} = \{1, 2, \ldots, N\}$：每架eVTOL为一个独立智能体，$N$在训练过程中由课程学习机制动态调整（从5逐步增至21）。
    \item \textbf{全局状态} $\mathcal{S}$：$S = (s_1, s_2, \ldots, s_N)$，其中$s_i$为第$i$架飞行器的局部观测。每个智能体仅能观测其自身状态与邻域（半径500m）内其他飞行器的信息，构成部分可观测设定。
    \item \textbf{联合动作} $\mathcal{A}$：$A = (a_1, a_2, \ldots, a_N)$，其中$a_i \in \mathbb{R}$为第$i$架飞行器的连续加速度指令。
    \item \textbf{状态转移} $\mathcal{P}(S' \mid S, A)$：受飞行器运动学、最大转弯角约束、风扰与传感器噪声的共同影响。
    \item \textbf{联合奖励} $\mathcal{R}$：$R = \frac{1}{N}\sum_{i=1}^{N} r_i$，所有飞行器奖励的均值。$r_i$由安全、速度、协同与效率四个分量加权构成（详见第3.3节）。
    \item \textbf{折扣因子} $\gamma$：取值为0.99。
\end{itemize}

优化目标为找到最优联合策略$\pi^* = (\pi_1^*, \pi_2^*, \ldots, \pi_N^*)$，最大化期望累计折扣回报：
\begin{equation}
    J(\pi) = \mathbb{E}_{\pi}\left[\sum_{t=0}^{T} \gamma^t R_t\right]
\end{equation}

鉴于部分可观测性，本文采用集中训练、分布式执行（CTDE）架构：训练阶段通过GraphMixer混合网络（第3.6节）利用全局状态信息学习协同策略；执行阶段每个智能体仅基于局部观测通过单智能体网络（第3.5节）做出独立决策，无需全局通信。

\subsection{环境模型}

本文设计的低空交通仿真环境\texttt{AirTrafficEnv}以二维平面模拟交叉路口空域（$1000\text{m} \times 1000\text{m}$），中心为坐标原点。飞行器从均匀分布于圆周的初始位置出发，方向指向中心。需要指出，当前环境是对低空交叉路口的简化抽象，在三维空间、多路口网络、复杂风场建模等维度上的扩展是未来工作的重要方向（详见第5.2节）。

\subsubsection{状态空间的三层定义}

为避免状态定义的歧义，本文将状态信息明确区分为三个层次：

\textbf{第一层——原始状态（5维）：}
每架飞行器$i$的原始状态向量为仅包含自身运动信息的5维向量：
\begin{equation}
    s_i^{\text{raw}} = [x_i,\; y_i,\; v_i,\; v_i^{\text{target}},\; d_i^{\text{center}}]
    \label{eq:raw_state}
\end{equation}
其中$(x_i, y_i)$为平面位置（单位：m），$v_i$为当前速度（单位：km/h），$v_i^{\text{target}}$为目标巡航速度（单位：km/h），$d_i^{\text{center}}$为到交叉中心的欧氏距离（单位：m）。

\textbf{第二层——局部观测（含邻域与历史信息）：}
在原始状态基础上，通过500m邻域半径与时间窗口$T=10$进行扩展，构建包含以下信息的局部观测张量：
\begin{itemize}
    \item 邻域飞行器$j \in \mathcal{N}_i$的相对位置$\Delta x_{ij}, \Delta y_{ij}$、相对速度$\Delta v_{ij}$、加速度$a_j$与航向角$\theta_j$；
    \item 时间窗口$T=10$内的历史轨迹序列（位置、速度、目标速度、到中心距离）。
\end{itemize}
该局部观测构成了单智能体网络的直接输入，维度随邻域大小动态变化。

\textbf{第三层——图节点特征（编码后256维）：}
局部观测经LSTM编码器（第3.5节）与DyGAT层（第3.4节）处理后，得到固定维度256维的时空耦合特征向量$h_i^{\text{coupled}}$，用于后续的策略输出与价值估计。三层信息处理的端到端流程为：原始状态$\to$局部观测$\to$图节点特征$\to$策略/价值输出。

\subsubsection{动作空间与物理约束}

动作空间为连续值，$a_i$对应第$i$架飞行器的加速度指令。为确保控制指令的物理可执行性，施加以下约束：
\begin{itemize}
    \item 最大加速度：$3.0\,\text{m/s}^2$，最大减速度：$4.0\,\text{m/s}^2$；
    \item 速度范围：$[60, 140]\,\text{km/h}$（与表\ref{tab:env_params}一致）；
    \item 最大转弯角：$\pi/6$（$30^\circ$），限制航向变化率以避免非物理机动。
\end{itemize}

动作施加模型为：$v_i^{t+1} = \text{clip}(v_i^{t} + a_i \cdot \Delta t \cdot 3.6,\; 60,\; 140)$，其中$\Delta t = 0.5\text{s}$，系数3.6完成m/s到km/h的单位转换。

\subsubsection{不确定性建模}

为模拟真实低空环境中的扰动，环境中引入两类不确定性：
\begin{enumerate}[label=(\arabic*)]
    \item \textbf{传感器噪声}：对位置观测添加$\mathcal{N}(0, 0.5^2)$的高斯噪声（单位：m），模拟GPS定位误差；
    \item \textbf{风扰}：为每架飞行器逐时间步添加随机方向$\theta_w \sim \mathcal{U}(0, 2\pi)$、固定强度$0.1$的风扰速度（单位：m/s）。
\end{enumerate}

当前风扰与噪声模型是对真实低空环境（如城市峡谷效应、Dryden风湍流谱）的高度简化抽象，其影响在第5.2节中进行了进一步讨论。

\subsubsection{动态安全距离与冲突判定}

传统固定安全距离无法适配不同飞行速度下的制动需求\cite{ref7}。本文采用与空域平均速度$\bar{v}$（单位：km/h）正相关的动态安全距离：
\begin{equation}
    d_{\text{safety}} = d_{\text{base}} + \bar{v} \times 0.5
    \label{eq:dynamicsafety}
\end{equation}
其中$d_{\text{base}} = 50\,\text{m}$为基础安全距离。当$\exists i \neq j,\; \|(x_i, y_i) - (x_j, y_j)\|_2 < d_{\text{safety}}$时，判定发生一次飞行冲突。

\subsection{分层奖励函数设计}

单架飞行器的即时奖励由四个分量加权构成：
\begin{equation}
    r_i = r_{\text{safe}} + r_{\text{speed}} + r_{\text{coop}} + r_{\text{eff}}
    \label{eq:total_reward}
\end{equation}

\subsubsection{安全核心奖励}

安全奖励采用指数惩罚机制对冲突风险进行强约束：
\begin{equation}
    r_{\text{safe}} =
    \begin{cases}
        -50 \cdot \exp\left(2 \cdot \dfrac{d_{\text{safety}} - d_{\min}}{d_{\text{safety}}}\right), & d_{\min} < d_{\text{safety}} \\[10pt]
        -10 \cdot \dfrac{1.2d_{\text{safety}} - d_{\min}}{0.2d_{\text{safety}}}, & d_{\text{safety}} \leq d_{\min} < 1.2d_{\text{safety}} \\[10pt]
        5, & d_{\min} \geq 1.2d_{\text{safety}}
    \end{cases}
    \label{eq:rsafe}
\end{equation}
其中$d_{\min} = \min_{j \neq i} \|(x_i, y_i) - (x_j, y_j)\|_2$。该设计具有三个特征：危险区（$d_{\min} < d_{\text{safety}}$）指数增长惩罚提供强安全约束；预警区（$d_{\text{safety}} \leq d_{\min} < 1.2d_{\text{safety}}$）线性过渡惩罚实现平滑梯度信号；安全区（$d_{\min} \geq 1.2d_{\text{safety}}$）固定正奖励维持稳定策略。

\subsubsection{速度控制奖励}

速度控制奖励的阈值设计与环境参数表\ref{tab:env_params}中的速度范围$[60, 140]\,\text{km/h}$保持一致：
\begin{equation}
    r_{\text{speed}} =
    \begin{cases}
        -50 \cdot \dfrac{v_i - 90}{50}, & v_i > 90\,\text{km/h} \\[8pt]
        -30 \cdot \dfrac{65 - v_i}{25}, & v_i < 65\,\text{km/h} \\[8pt]
        20, & 65 \leq v_i \leq 90\,\text{km/h}
    \end{cases}
    \label{eq:rspeed}
\end{equation}
最优速度区间$[65, 90]\,\text{km/h}$位于飞行包线中低段，鼓励飞行器以安全适中的速度飞行。

\subsubsection{协同增效奖励与运行效率奖励}

协同奖励鼓励邻域速度一致性：
\begin{equation}
    r_{\text{coop}} = -0.2 \cdot |v_i - \bar{v}_{\text{neighbor}}|
    \label{eq:rcoop}
\end{equation}
其中$\bar{v}_{\text{neighbor}}$为邻域（半径500m）内飞行器的平均速度。

效率奖励惩罚与目标速度的偏差：
\begin{equation}
    r_{\text{eff}} = -0.5 \cdot |v_i - v_i^{\text{target}}|
    \label{eq:reff}
\end{equation}

全局奖励为所有飞行器奖励的均值。当所有飞行器均进入目标区域（距中心100m以内）时，额外给予$+200$的全局完成奖励作为稀疏终端信号。

\subsection{动态图注意力网络 (DyGAT)}

传统GAT仅能处理静态图结构，无法刻画飞行器间随运动不断变化的时空关联。本文设计的DyGAT同时融合空间注意力（同时间步邻域关系）、时序注意力（历史轨迹演化趋势）与动态边权重（基于距离-速度-航向角的物理耦合强度），三者协同实现对飞行器耦合关系的精准刻画。

\subsubsection{空间注意力机制}

对于节点特征矩阵$h \in \mathbb{R}^{N \times F_{\text{in}}}$（$N$为当前飞行器数量，$F_{\text{in}}$为输入特征维度），首先经线性变换$W$得到特征映射，随后计算注意力分数：
\begin{equation}
    e_{ij} = \text{LeakyReLU}_{0.2}\left(a^{\top} \cdot [Wh_i \parallel Wh_j]\right)
    \label{eq:spatial_attn}
\end{equation}
其中$\parallel$为特征拼接，$a \in \mathbb{R}^{2F_{\text{out}}}$为可学习注意力权重向量。经softmax归一化（仅对$adj_{ij} > 0$的邻域节点）：
\begin{equation}
    \alpha_{ij} = \frac{\exp(e_{ij})}{\sum_{k \in \mathcal{N}_i} \exp(e_{ik})}
    \label{eq:alpha}
\end{equation}
空间注意力的输出$h_i^{\text{spatial}} = \sum_{j \in \mathcal{N}_i} \alpha_{ij} \cdot Wh_j$实现了邻域特征的加权聚合。

\subsubsection{时序注意力机制}

时序注意力融合历史$T$步轨迹信息。对于历史特征序列$h_{\text{history}} \in \mathbb{R}^{T \times N \times F_{\text{out}}}$：
\begin{align}
    e_{t,i} &= \text{LeakyReLU}_{0.2}\left(a_t^{\top} \cdot [Wh_i^{\text{curr}} \parallel Wh_{t,i}^{\text{hist}}]\right) \label{eq:temporal_attn} \\
    \beta_{t,i} &= \frac{\exp(e_{t,i})}{\sum_{\tau=0}^{T-1} \exp(e_{\tau,i})} \label{eq:beta} \\
    h_i^{\text{temporal}} &= \sum_{t=0}^{T-1} \beta_{t,i} \cdot Wh_{t,i}^{\text{hist}} \label{eq:temporal_out}
\end{align}
最终时空融合特征为$h'_i = h_i^{\text{spatial}} + 0.3 \cdot h_i^{\text{temporal}}$，时序信息权重$0.3$经实验调优确定。

\subsubsection{动态边权重计算}

为替代静态二值邻接矩阵，基于物理量计算连续动态边权重：
\begin{equation}
    w_{ij} = \underbrace{\exp(-d_{ij}/1000)}_{\text{距离权重}} \cdot \underbrace{\exp(-|v_i - v_j|/20)}_{\text{速度权重}} \cdot \underbrace{\exp(-|\theta_i - \theta_j|/(\pi/4))}_{\text{航向权重}}
    \label{eq:edge_weight}
\end{equation}
三个指数核分别在距离、速度差和航向差维度上度量飞行器对的物理耦合强度。距离越近、速度越接近、航向越一致的飞行器对获得更高权重——模型据此将有限注意力资源聚焦于高风险交互对。动态边权重矩阵$W = [w_{ij}]$与静态邻接矩阵$adj$通过逐元素乘（Hadamard积）融合。

\subsection{单智能体网络}

单智能体网络为每架eVTOL提供独立的策略$\pi_i(a_i \mid o_i)$与价值估计$Q_i(o_i, a_i)$：

\begin{enumerate}[label=(\arabic*)]
    \item \textbf{LSTM编码器}：2层LSTM（隐藏维度64，Dropout 0.1）对$T=10$步历史轨迹序列进行编码，输出$h_i^{\text{traj}} \in \mathbb{R}^{64}$作为轨迹特征；
    \item \textbf{DyGAT层}：以$h_i^{\text{traj}}$为节点特征，通过DyGAT（第3.4节）聚合邻域信息，输出$h_i^{\text{coupled}} \in \mathbb{R}^{256}$；
    \item \textbf{策略头}：2层MLP（256$\to$256$\to$1，Tanh激活）输出动作均值$\mu_i$。动作从$\mathcal{N}(\mu_i, 0.1^2)$采样，训练时标准差固定为0.1提供探索噪声；
    \item \textbf{Q值头}：2层MLP（256$\to$256$\to$128$\to$1，ReLU激活）输出状态-动作价值$Q_i$。
\end{enumerate}

\subsection{GraphMixer混合网络}

为实现CTDE架构下的全局协同，本文采用GraphMixer替代传统QMIX的超网络混合结构\cite{ref11}。GraphMixer将各智能体4维状态向量$(x_i, y_i, v_i, v_i^{\text{target}})$作为图节点特征，在完全图上通过GAT层（第3.4.1节中的空间注意力机制）进行信息聚合，取节点均值作为全局表征$h_{\text{global}} \in \mathbb{R}^{256}$，再经MLP（256$\to$256$\to$1）输出标量全局Q值$Q_{\text{tot}}$。

GraphMixer相较QMIX具有两个关键优势：(1)~GAT的注意力权重可自适应学习智能体间的非线性耦合关系，而非QMIX中受限的单调性约束；(2)~完全图结构天然适配飞行器数量$N$动态变化的课程学习场景，无需重新设计超网络输入维度。

\subsection{训练优化机制}

\subsubsection{课程学习管理器}

课程学习机制将训练划分为递进阶段。初始阶段：$N=5$，$d_{\text{base}}=80\text{m}$。进阶条件：连续3次综合性能评估达标（冲突率$<$阈值且通行完成率$>80\%$），则$N \leftarrow \min(N+2,\; 21)$，$d_{\text{base}} \leftarrow \max(d_{\text{base}}-5,\; 50)$。回退机制：连续5次性能下降时，$N \leftarrow \max(N-1,\; 5)$。此设计确保模型先掌握基础冲突规避策略，再逐步适应高密度场景。

\subsubsection{基于冲突事件的轨迹加权策略}

PPO通过裁剪代理目标函数保证策略更新的稳定性\cite{ref11,ref17}：
\begin{equation}
    L^{\text{CLIP}}(\theta) = \hat{\mathbb{E}}_t\left[\min\left(r_t(\theta)\hat{A}_t,\; \text{clip}(r_t(\theta), 1-\epsilon, 1+\epsilon)\hat{A}_t\right)\right]
    \label{eq:ppo_clip}
\end{equation}
其中$r_t(\theta) = \pi_\theta(a_t \mid s_t) / \pi_{\theta_{\text{old}}}(a_t \mid s_t)$为策略概率比，$\epsilon = 0.2$，$\hat{A}_t$为广义优势估计（GAE）。

针对冲突样本高度稀疏的问题，本文在\textbf{保持PPO on-policy理论框架不变}的前提下，对episode级训练轨迹施加差异化损失权重：
\begin{equation}
    L_{\text{total}} = \lambda \cdot \left(L_{\text{policy}} + L_{q} + 0.5 \cdot L_{\text{value}}\right)
    \label{eq:weighted_loss}
\end{equation}
其中$\lambda = 2.0$（若episode内发生冲突事件），$\lambda = 1.0$（若无冲突）。该设计的直觉是：包含冲突事件的episode是模型最需要从中学习的"高价值轨迹"——通过在其上施加更大的梯度更新步长，等效于提高了稀疏冲突样本的利用效率，同时避免了引入off-policy机制破坏PPO的理论保证。

训练中采用梯度裁剪（$\max\_\text{norm}=1.0$）与指数学习率衰减（$\gamma_{\text{lr}}=0.99$，初始学习率$5 \times 10^{-4}$）。

\subsection{模型整体架构}

整体架构如图\ref{fig:architecture}所示，由四个层次构成：
\begin{enumerate}[label=(\arabic*)]
    \item \textbf{环境层}：\texttt{AirTrafficEnv}提供融合风扰与传感器噪声的多飞行器状态；
    \item \textbf{编码层}：LSTM编码器提取轨迹特征，DyGAT层实现时空耦合聚合；
    \item \textbf{决策层}：策略头输出加速度均值，Q值头输出状态-动作价值，GraphMixer输出全局Q值；
    \item \textbf{训练层}：课程学习控制难度递进，冲突轨迹加权提升稀疏样本效率，PPO保证更新稳定性。
\end{enumerate}

\begin{figure}[H]
    \centering
    \includegraphics[width=1.0\textwidth]{QQ图片20260523201329.png}
    \caption{模型整体架构}
    \label{fig:architecture}
\end{figure}

%==========================================================================
\section{实验与结果分析}
%==========================================================================

\subsection{实验设置}

\subsubsection{仿真环境参数}

实验基于PyTorch 2.0实现，硬件为NVIDIA RTX 4060 GPU与Intel i9-14900K CPU，软件为Python 3.14。仿真核心参数如表\ref{tab:env_params}所示。

\begin{table}[H]
    \centering
    \caption{仿真环境核心参数}
    \label{tab:env_params}
    \begin{tabular}{|c|c|}
        \hline
        \textbf{参数名称} & \textbf{参数值} \\
        \hline
        交叉空域大小 & $1000\text{m} \times 1000\text{m}$ \\
        基础安全距离 & $50\text{m}$ \\
        时间步长 $\Delta t$ & $0.5\text{s}$ \\
        最大仿真步长 & $200$ \\
        飞行器目标速度 & $70 \pm 10$ km/h \\
        速度范围 & $[60, 140]$ km/h \\
        最大加速度 & $3.0$ m/s$^2$ \\
        最大减速度 & $4.0$ m/s$^2$ \\
        最大转弯角 & $\pi/6$ rad \\
        时间窗口 $T$ & $10$ \\
        邻域半径 & $500\text{m}$ \\
        训练轮次 & $1000$ \\
        折扣因子 $\gamma$ & $0.99$ \\
        PPO裁剪参数 $\epsilon$ & $0.2$ \\
        初始学习率 & $5 \times 10^{-4}$ \\
        \hline
    \end{tabular}
\end{table}

\subsubsection{基线方法}

选取8种基线方法，按类别划分为四组：

\begin{enumerate}[label=\textbf{第\arabic*组：}]
    \item \textbf{传统规则方法}
    \begin{itemize}
        \item \textbf{RR（Round Robin）}：固定轮询调度，各方向飞行器按固定顺序获得通行时隙；
        \item \textbf{FCFS（First-Come-First-Served）}：按到达交叉路口的时间顺序分配通行权。
    \end{itemize}
    \item \textbf{单智能体强化学习}
    \begin{itemize}
        \item \textbf{传统DQN}：集中式架构，全局状态输入，输出所有飞行器的速度控制指令。
    \end{itemize}
    \item \textbf{独立多智能体学习}
    \begin{itemize}
        \item \textbf{独立A2C}：每架飞行器为独立A2C智能体，无智能体间显式信息交互。
    \end{itemize}
    \item \textbf{集中训练分布式执行MARL（CTDE-MARL）}
    \begin{itemize}
        \item \textbf{MADDPG}：集中式Critic + 分布式Actor架构；
        \item \textbf{COMA}：反事实基线解决多智能体信用分配；
        \item \textbf{QMIX}：超网络混合结构的价值分解方法；
        \item \textbf{MAPPO}：多智能体PPO，目前MARL领域广泛使用的基准。
    \end{itemize}
\end{enumerate}

\subsubsection{评价指标}

选取4个核心评价指标：
\begin{enumerate}[label=(\arabic*)]
    \item \textbf{累计奖励}（Cumulative Reward）：每episode的全局奖励总和$\sum_t R_t$；
    \item \textbf{总冲突次数}（Total Conflicts）：每episode中所有时间步冲突对数之和（同一飞行器对在同一时间步的冲突计为1）；
    \item \textbf{平均延迟}（Average Delay）：所有飞行器因速度偏离目标速度产生的运行延迟均值；
    \item \textbf{通行完成率}（Passage Completion Rate）：成功到达目标区域（距中心100m内）的飞行器比例。
\end{enumerate}

\subsection{训练收敛性分析}

模型训练1000轮的收敛曲线如图\ref{fig:training_curve}所示。前200轮为快速收敛期：累计奖励从初始的大幅负值持续上升，冲突次数与平均延迟快速下降。200轮后进入稳定收敛期：奖励波动幅度$<5\%$，冲突次数稳定在低位。课程学习在约第50、120、200、300轮触发难度进阶（$5\to7\to9\to11\to13\to15\to17\to19\to21$架），每次进阶后出现短暂性能回落（1-3轮），随后快速恢复并超越进阶前水平——验证了课程学习机制在保证训练稳定性方面的关键作用。

\begin{figure}[H]
    \centering
    \includegraphics[width=0.8\textwidth]{QQ图片20260504231524.jpg}
    \caption{训练收敛曲线（累计奖励、冲突次数、平均延迟随训练轮次的变化）}
    \label{fig:training_curve}
\end{figure}

\subsection{不同方法的性能对比}

21架eVTOL测试场景下的综合性能对比如表\ref{tab:comparison}所示。

\begin{table}[H]
    \centering
    \caption{不同方法的性能对比（21架eVTOL，5次独立测试的均值$\pm$标准差）}
    \label{tab:comparison}
    \begin{tabular}{|c|c|c|c|c|}
        \hline
        \textbf{方法（类别）} & \textbf{累计奖励$\uparrow$} & \textbf{冲突次数$\downarrow$} & \textbf{平均延迟(s)$\downarrow$} & \textbf{通行完成率(\%)$\uparrow$} \\
        \hline
        RR（规则） & $-1256.3 \pm 89.2$ & $42.6 \pm 5.3$ & $1.58 \pm 0.12$ & $89.2 \pm 3.5$ \\
        FCFS（规则） & $-1187.5 \pm 76.4$ & $38.9 \pm 4.7$ & $1.42 \pm 0.10$ & $91.5 \pm 2.8$ \\
        传统DQN（单RL） & $-892.5 \pm 65.3$ & $28.3 \pm 3.9$ & $1.21 \pm 0.09$ & $94.7 \pm 2.1$ \\
        独立A2C（独立ML） & $-657.8 \pm 52.6$ & $19.7 \pm 2.8$ & $0.93 \pm 0.07$ & $96.3 \pm 1.7$ \\
        MADDPG（CTDE） & $-215.6 \pm 41.2$ & $25.4 \pm 3.2$ & $0.87 \pm 0.06$ & $95.8 \pm 1.9$ \\
        COMA（CTDE） & $187.3 \pm 35.7$ & $22.1 \pm 2.9$ & $0.79 \pm 0.06$ & $97.1 \pm 1.5$ \\
        QMIX（CTDE） & $426.5 \pm 28.9$ & $23.1 \pm 2.7$ & $0.76 \pm 0.05$ & $97.5 \pm 1.3$ \\
        MAPPO（CTDE） & $892.4 \pm 23.5$ & $20.3 \pm 2.4$ & $0.65 \pm 0.04$ & $98.2 \pm 1.1$ \\
        \textbf{本文方法} & $\mathbf{1286.7 \pm 18.2}$ & $\mathbf{18.2 \pm 2.1}$ & $\mathbf{0.57 \pm 0.03}$ & $\mathbf{99.1 \pm 0.8}$ \\
        \hline
    \end{tabular}
\end{table}

从表\ref{tab:comparison}可得出以下结论：

\textbf{(1) 与传统规则方法对比：} 相比RR，冲突次数降低57.2\%，延迟缩短63.9\%，通行完成率提升9.9个百分点；相比FCFS，冲突次数降低53.2\%，延迟缩短59.9\%。这一量级的提升体现了学习型方法相比固定规则的本质优势。

\textbf{(2) 与当前最优CTDE-MARL方法对比：} 相比MAPPO，冲突次数降低10.3\%（20.3$\to$18.2），延迟缩短12.3\%（0.65s$\to$0.57s），累计奖励提升44.2\%（892.4$\to$1286.7）。DyGAT的时空耦合建模、课程学习的渐进式训练与冲突轨迹加权机制的叠加效应共同贡献了这一性能增益。

\textbf{(3) 跨类别性能梯度：} 规则方法性能最差但计算成本最低（$<$1ms/步）；DQN在飞行器数量增多时面临状态维度爆炸；独立A2C由于缺乏智能体间协同导致策略收敛于局部最优；CTDE-MARL方法中MAPPO最优，本文方法在此基础上进一步提升。

\textbf{(4) QMIX的超网络瓶颈：} QMIX的超网络结构对输入维度有严格要求，在课程学习中飞行器数量$N$动态变化时需重新设计网络，性能受限（冲突次数23.1 $\pm$ 2.7）。GraphMixer通过完全图上的GAT聚合规避了此问题。

\subsection{消融实验}

表\ref{tab:ablation}报告了消融实验结果，量化了各核心模块的独立性能贡献。

\begin{table}[H]
    \centering
    \caption{消融实验结果（21架eVTOL，5次独立测试）}
    \label{tab:ablation}
    \begin{tabular}{|c|c|c|c|c|}
        \hline
        \textbf{模型配置} & \textbf{累计奖励$\uparrow$} & \textbf{冲突次数$\downarrow$} & \textbf{平均延迟(s)$\downarrow$} & \textbf{通行完成率(\%)$\uparrow$} \\
        \hline
        完整模型（本文） & $\mathbf{1286.7 \pm 18.2}$ & $\mathbf{18.2 \pm 2.1}$ & $\mathbf{0.57 \pm 0.03}$ & $\mathbf{99.1 \pm 0.8}$ \\
        $-$ DyGAT（改用标准GAT） & $215.3 \pm 32.6$ & $49.8 \pm 4.5$ & $0.79 \pm 0.05$ & $94.3 \pm 1.9$ \\
        $-$ 课程学习（直接21架训练） & $587.2 \pm 27.4$ & $28.4 \pm 3.1$ & $0.68 \pm 0.04$ & $97.2 \pm 1.4$ \\
        $-$ 冲突轨迹加权（$\lambda\equiv1$） & $763.5 \pm 24.1$ & $24.6 \pm 2.8$ & $0.65 \pm 0.04$ & $97.8 \pm 1.2$ \\
        $-$ 课程学习 $-$ 冲突轨迹加权 & $128.6 \pm 38.5$ & $35.7 \pm 3.9$ & $0.82 \pm 0.06$ & $95.6 \pm 1.7$ \\
        \hline
    \end{tabular}
\end{table}

消融实验的核心发现如下：

\textbf{DyGAT是性能贡献最大的单一模块：} 替换为标准GAT后，冲突次数从18.2激增至49.8（上升173.6\%），通行完成率从99.1\%降至94.3\%。这一幅度远超其他模块的独立贡献，验证了动态边权重与时空注意力融合是模型实现精准冲突规避的根本机制。

\textbf{课程学习与冲突轨迹加权形成互补增效：} 单独移除课程学习导致冲突次数上升56.0\%（18.2$\to$28.4），单独移除冲突轨迹加权导致冲突次数上升35.2\%（18.2$\to$24.6）。同时移除两者时冲突次数上升96.2\%（18.2$\to$35.7），表明两种机制存在协同增效——课程学习提供从简到繁的渐进训练路径，冲突轨迹加权确保在每个难度等级中关键样本被充分学习。

\subsection{不同飞行器密度下的鲁棒性分析}

在$N \in \{5, 10, 15, 21\}$四种密度下的冲突次数对比如图\ref{fig:density_robustness}所示。

\begin{figure}[H]
    \centering
    \includegraphics[width=0.65\textwidth]{QQ图片20260524171218.png}
    \caption{不同飞行器密度下的冲突次数对比}
    \label{fig:density_robustness}
\end{figure}

随着$N$从5增至21，本文方法的冲突次数增长斜率（$\approx 0.75$/架）显著低于MAPPO（$\approx 1.15$/架）与RR（$\approx 2.4$/架）。在低密度（$N=5$）场景中，所有CTDE-MARL方法的冲突次数均较低（本文2.1次，MAPPO 2.4次，QMIX 2.8次），差异不显著；密度增至$N=15$时，本文方法相比MAPPO的冲突次数优势扩大至22\%；在$N=21$的高密度极值下，本文方法的冲突次数仍控制在19次以内，而RR已突破40次。这一趋势表明DyGAT的时空耦合建模能力在高密度场景中优势愈发显著——飞行器越多，精准区分不同邻域飞行器的冲突风险等级越关键。

\subsection{敏感性分析}

在$N=21$场景下，分别改变基础安全距离$d_{\text{base}}$（40--70m）、风扰强度（0--0.3）与传感器噪声标准差$\sigma$（0--1.5m），测试模型性能的变化幅度，结果如图\ref{fig:sensitivity_analysis}所示。

\begin{figure}[H]
    \centering
    \includegraphics[width=0.8\textwidth]{QQ图片20260524162300.png}
    \caption{敏感性分析结果}
    \label{fig:sensitivity_analysis}
\end{figure}

\textbf{安全距离参数：} $d_{\text{base}}$从40m增至70m时，冲突次数从22.5降至14.7（$\downarrow$34.7\%），延迟从0.52s升至0.63s（$\uparrow$21.2\%），符合安全-效率的基本权衡关系。模型能够在不同安全要求下调整策略，在安全与效率之间取得合理平衡。

\textbf{风扰强度：} 风扰从0增至0.3时，冲突次数从16.3升至20.8（$\uparrow$27.6\%），性能退化幅度温和。DyGAT的时序注意力机制在此发挥了关键作用——通过学习历史轨迹的趋势信息，模型能够部分补偿风扰带来的位置不确定性。

\textbf{传感器噪声：} $\sigma$从0增至1.5m时，冲突次数从17.1升至19.6（$\uparrow$14.6\%），是所有参数中影响最小的维度。模型对传感器噪声的较高鲁棒性得益于LSTM编码器对时序状态序列的平滑作用。

\subsection{计算成本分析}

图\ref{fig:computation_cost}展示了不同$N$下各方法的单步推理耗时。本文方法在$N=21$时的平均推理时间为$4.5\text{ms}$，远低于仿真时间步长$500\text{ms}$（实时性裕度$\approx 111\times$）。计算时间从$N=5$时的$1.2\text{ms}$到$N=21$时的$4.5\text{ms}$，增长斜率（$\approx 0.21\text{ms}$/架）在所有CTDE-MARL方法中最低。COMA在$N=21$时耗时超过$15\text{ms}$，QMIX为$2.1\text{ms}$（$N=5$）到约$8\text{ms}$（$N=21$）。

\begin{figure}[H]
    \centering
    \includegraphics[width=0.8\textwidth]{QQ图片20260524163903.jpg}
    \caption{不同飞行器数量下的平均单步推理计算时间对比}
    \label{fig:computation_cost}
\end{figure}

%==========================================================================
\section{讨论与结论}
%==========================================================================

\subsection{结果解读与贡献总结}

本文提出的耦合多智能体强化学习管控模型在低空交叉路口场景中取得的性能优势，可归因于三个核心设计要素的协同作用：

\textbf{(1) DyGAT的精准时空耦合刻画是性能基石。} 消融实验中DyGAT的移除导致冲突次数增加173.6\%，在所有消融项中贡献最大。其机制可解释为：动态边权重（式\ref{eq:edge_weight}）使模型能够根据距离、速度差与航向差三个物理量自动区分邻域飞行器的风险等级——距离近且航向冲突的飞行器对被赋予最高注意力权重，模型据此优先规避最紧迫的冲突威胁。

\textbf{(2) 课程学习与冲突轨迹加权的协同效应。} 课程学习（5$\to$21架渐进训练）确保模型在每个难度阶段都能充分收敛后再进阶，避免了直接高密度训练的策略崩溃。冲突轨迹加权（$\lambda=2.0$）在保持PPO on-policy理论纯净性的前提下，有效利用了珍贵的冲突事件训练信号。两者协同——课程学习铺就平滑训练路径，轨迹加权使每段路径上的关键样本被高效利用。

\textbf{(3) GraphMixer的灵活全局协同。} 相比QMIX的超网络对固定输入维度的依赖，GraphMixer的完全图GAT聚合天然适配$N$变化的课程学习场景，同时注意力权重可学习非单调的智能体间价值关系。

\subsection{方法局限性}

本文方法存在以下局限性，需要在未来工作中解决：

\textbf{(1) 场景维度限制：} 当前环境为二维单交叉路口简化模型。真实低空交通涉及三维高度层（含垂直起降阶段）、多路口网络拓扑、空域扇区划分以及起降场容量约束等复杂因素。

\textbf{(2) 理想通信假设：} 当前CTDE架构的训练阶段假设全局状态可无延迟获取。在真实UTM部署中，通信延迟、数据包丢失与有限带宽将影响集中式训练的质量与分布式执行的协同效率。

\textbf{(3) 飞行器同质化假设：} 所有飞行器使用相同的动力学参数（最大加速度$3.0\text{m/s}^2$等）。实际运营中不同制造商与不同载荷等级的eVTOL具有差异化的飞行包线。

\textbf{(4) 扰动建模的简化：} 风扰模型为均匀随机分布，噪声模型为高斯白噪声。真实低空风场具有时空相关性（如城市峡谷效应的非均匀风场），传感器误差也具有时间偏置特性。

\textbf{(5) 基线比较范围：} 本文未与模型预测控制（MPC）、速度障碍法（ORCA/RVO）与控制障碍函数（CBF）等航空安全关键系统的标准方法进行对比。这些方法在提供形式化安全保证方面具有独特优势，与它们的全面比较将更准确地定位本文方法的性能边界。

\subsection{实际应用潜力}

在充分认识上述局限性的前提下，本文模型具有以下应用潜力：(1)~可作为低空UTM冲突管理模块的候选算法，为高密度eVTOL自主协同管控提供技术储备；(2)~分布式执行架构使得计算负载分散至各飞行器机载设备，避免了集中式决策的单点故障风险；(3)~模型对风扰与传感器噪声表现出的鲁棒性（性能退化$<$28\%），表明其在非理想感知条件下仍具备一定的可靠性。在进一步解决通信约束、异构动力学与实飞验证后，该技术路线具备向工程化部署过渡的潜力。

\subsection{未来工作}

面向工程化部署的实际需求，未来研究将从以下五个方向展开：

\begin{enumerate}[label=(\arabic*)]
    \item \textbf{三维广域空域扩展：} 将模型迁移至三维空间，引入高度层分配、垂直起降动力学与多路口网络拓扑，研究区域级交通流协同优化。
    \item \textbf{通信受限条件下的鲁棒训练：} 在CTDE框架中显式建模通信延迟与丢包，探索联邦学习或延迟容忍的信息共享机制以保证通信退化条件下的协同质量。
    \item \textbf{与安全关键方法的系统对比：} 引入MPC、ORCA与CBF等基线方法，在当前仿真环境中建立全面的性能基准。
    \item \textbf{异构飞行器群体建模：} 支持至少3类不同动力学参数的eVTOL混合运行，验证模型在异构场景中的泛化性。
    \item \textbf{硬件在环仿真与实飞验证：} 结合真实飞控硬件与飞行日志数据，开展硬件在环测试，推动模型从仿真验证向工程落地的关键跨越。
\end{enumerate}

%==========================================================================
% 参考文献（按正文首次引用顺序排列）
%==========================================================================
\begin{thebibliography}{99}

% [1] 首次引用：第1节——Agogino & Tumer耦合智能体调控框架
\bibitem{ref1} AGOGINO A, TUMER K. Regulating air traffic flow with coupled agents[C]//Proceedings of the 7th International Conference on Autonomous Agents and Multiagent Systems (AAMAS 2008). Richland: IFAAMAS, 2008: 535-542.

% [2] 首次引用：第1节——Agogino & Tumer多智能体空中交通流量管理
\bibitem{ref2} AGOGINO A K, TUMER K. A multiagent approach to managing air traffic flow[J]. Autonomous Agents and Multi-Agent Systems, 2012, 24(1): 1-25.

% [3] 首次引用：第1节——Brittain & Wei自主间隔保持
\bibitem{ref3} BRITTAIN M, WEI P. Autonomous separation assurance in a high-density en route sector: a deep multi-agent reinforcement learning approach[C]//2019 IEEE Intelligent Transportation Systems Conference (ITSC). Piscataway: IEEE, 2019: 3256-3262.

% [4] 首次引用：第1节——Tumer & Agogino分布式智能体流量管理
\bibitem{ref4} TUMER K, AGOGINO A. Distributed agent-based air traffic flow management[C]//Proceedings of the 6th International Joint Conference on Autonomous Agents and Multiagent Systems (AAMAS 2007). Honolulu: ACM, 2007: 330-337.

% [5] 首次引用：第1节——Tumer & Agogino学习型多智能体ATM
\bibitem{ref5} TUMER K, AGOGINO A. Improving air traffic management with a learning multiagent system[J]. IEEE Intelligent Systems, 2009, 24(1): 18-21.

% [6] 首次引用：第1节——Xie等RL流量管理
\bibitem{ref6} XIE Y B, GARDI A, SABATINI R. Reinforcement learning-based flow management techniques for urban air mobility and dense low-altitude air traffic operations[C]//2021 IEEE/AIAA 40th Digital Avionics Systems Conference (DASC). Piscataway: IEEE, 2021: 1-10.

% [7] 首次引用：第1节——Spatharis等层级MARL空管方案
\bibitem{ref7} SPATHARIS C, BASTAS A, KRAVARIS T, et al. Hierarchical multiagent reinforcement learning schemes for air traffic management[J]. Neural Computing and Applications, 2021, 33(14): 8649-8671.

% [8] 首次引用：第1节——张春阳AI交通控制综述
\bibitem{ref8} 张春阳, 席雅琪. 人工智能在城市交通控制中的应用综述[J]. 信息系统工程, 2024(07): 33-36.

% [9] 首次引用：第1节——Ghosh等深度集成MARL
\bibitem{ref9} GHOSH S, LAGUNA S, LIM S H, et al. A deep ensemble method for multi-agent reinforcement learning: a case study on air traffic control[C]//Proceedings of the 31st International Conference on Automated Planning and Scheduling (ICAPS 2021). Palo Alto: AAAI Press, 2021: 468-476.

% [10] 首次引用：第1节——Aslani等Actor-Critic交通信号
\bibitem{ref10} ASLANI M, MESGARI M S, WIERING M. Adaptive traffic signal control with actor-critic methods in a real-world traffic network with different traffic disruption events[J]. Transportation Research Part C: Emerging Technologies, 2017, 85: 732-752.

% [11] 首次引用：第1节——翟子洋等RL/DRL交通信号综述
\bibitem{ref11} 翟子洋, 郝茹茹, 董世浩. 大规模智慧交通信号控制中的强化学习和深度强化学习方法综述[J]. 计算机应用研究, 2024, 41(6): 1618-1627.

% [12] 首次引用：第1节——Ghavamzadeh层级MARL
\bibitem{ref12} GHAVAMZADEH M, MAHADEVAN S, MAKAR R. Hierarchical multi-agent reinforcement learning[J]. Autonomous Agents and Multi-Agent Systems, 2006, 13(2): 197-229.

% [13] 首次引用：第2节——王迎城市交通DRL硕士论文
\bibitem{ref13} 王迎. 城市交通信号控制深度强化学习算法研究[D]. 济南: 山东交通学院, 2024.

% [14] 首次引用：第2节——赵晓华Q-learning+BP交叉口控制
\bibitem{ref14} 赵晓华, 石建军, 李振龙, 赵国勇. 基于Q-learning和BP神经元网络的交叉口信号灯控制[J]. 公路交通科技, 2007, 24(7): 99-102.

% [15] 首次引用：第2节——承向军Q-学习交通信号控制
\bibitem{ref15} 承向军, 常歆识, 杨肇夏. 基于Q-学习的交通信号控制方法[J]. 系统工程理论与实践, 2006(8): 136-140.

% [16] 首次引用：第2节——赖建辉DRL交通控制硕士论文
\bibitem{ref16} 赖建辉. 基于深度强化学习的交通控制优化方法研究及实现[D]. 杭州: 浙江工业大学, 2020.

% [17] 首次引用：第2节——张新武DRL交通信号优化硕士论文
\bibitem{ref17} 张新武. 基于深度强化学习算法的交通信号控制优化研究[D]. 太原: 太原科技大学, 2024.

% [18] 首次引用：第2节——卢守峰在线Q学习交叉口模型
\bibitem{ref18} 卢守峰, 张术, 刘喜敏. 平均排队长度差最小的单交叉口在线Q学习模型[J]. 公路交通科技, 2014, 31(11): 116-122.

% [19] 首次引用：第2节——江波DQN机场道口调速
\bibitem{ref19} 江波, 李彦冬, 刘威陇, 等. 应用深度Q学习的机场道口协同智能调速方法[J]. 航空计算技术, 2022, 52(1): 26-30.

\end{thebibliography}

\end{document}
